%!TEX root = ../main.tex

\chapter{Introduction\label{chap:introduction}}
Precise time synchronization is an important requirement in distributed UAV systems, where sensor data, control actions, and diagnostic messages must be aligned across multiple nodes. \cite{r1,r2,r7} This thesis focuses on the design and validation of a compact STM32-based synchronization module that combines DCF77 absolute time reception with RTK PPS-based second generation and ROS2-based evaluation.

The chapter introduces the motivation for the work, defines the main objectives, and summarizes the structure of the thesis.

\section{Motivation}

The increasing use of UAV swarms in applications such as cooperative mapping, formation flight, and distributed sensing places high demands on consistent timing between individual nodes. \cite{r1,r2,r7} In such systems, sensor measurements, control actions, and diagnostic messages must be timestamped in a comparable time base so that data from multiple vehicles can be correctly correlated.

Many synchronization approaches rely on communication networks or external computing resources. However, wireless communication in UAV systems can introduce variable latency, packet delay, and jitter, especially in dynamic multi-robot environments. \cite{lopez2021uavmesh} These effects make purely network-based synchronization unsuitable as the only timing source when deterministic timing behavior is required.

This thesis is motivated by the need for a compact embedded synchronization module that can provide reliable local time information independently of network delay. The developed solution uses a DCF77 radio time signal to obtain absolute calendar time and an RTK receiver PPS signal to generate a stable one-second reference. An STM32 microcontroller decodes and validates the time information, maintains a software clock, and transmits synchronization diagnostics to a ROS2-based evaluation environment.

The goal is therefore not only to design a theoretical synchronization concept, but also to implement and experimentally validate a practical hardware and firmware pipeline suitable for future UAV swarm timing experiments.

\section{Objectives}

The main objective of this thesis is to design, implement, and evaluate a microcontroller-based time synchronization module for UAV swarm-related experiments. The work focuses on creating a compact embedded system capable of acquiring absolute time, maintaining a local software clock, and validating synchronization behavior through a ROS2-based evaluation pipeline.
\\
The specific objectives are:
\begin{itemize}
    \item to analyze time synchronization principles relevant to distributed UAV and multi-robot systems,
    \item to design a custom STM32-based hardware platform with interfaces for DCF77, RTK PPS, UART communication, and debugging,
    \item to implement firmware for DCF77 frame decoding, parity validation, software clock maintenance, and RTK PPS-driven second generation,
    \item to integrate the synchronization module with a companion-computer environment using UART and ROS2,
    \item to experimentally evaluate synchronization stability, communication reliability, and timing behavior using logged measurement data.
\end{itemize}

The main contribution of this thesis is the design and validation of a complete embedded synchronization prototype. The work includes the custom STM32-based PCB, firmware for DCF77 decoding and RTK PPS-driven clock generation, UART diagnostic communication, and a ROS2-based logging and analysis pipeline. The validated result is stable DCF77-based absolute time acquisition and PPS-driven second generation with no observable long-term software-clock drift at millisecond-level measurement resolution. The work does not claim sub-microsecond synchronization accuracy, because the final evaluation was performed using firmware diagnostics, UART messages, ROS2 logs, and Python post-processing rather than direct hardware-level timestamping.

\section{Structure of the Thesis}

The thesis is structured into eight chapters. Chapter 1 introduces the motivation, objectives, and scope of the work. Chapter 2 presents the theoretical background of time synchronization, clock drift, GNSS/RTK timing, DCF77 time reception, and microcontroller-based timing architectures. Chapter 3 describes the overall system design and explains the role of the main synchronization components.

Chapter 4 focuses on the hardware implementation of the custom STM32-based PCB, including the power supply, signal interfaces, PCB layout, and communication paths. Chapter 5 describes the firmware implementation, including DCF77 decoding, parity validation, software clock maintenance, RTK PPS-based second generation, and UART diagnostic output.

Chapter 6 presents the integration of the synchronization module with the ROS2 environment used for data collection and analysis. Chapter 7 describes the measurement setup, logged results, synchronization behavior, and practical limitations of the evaluation method. Finally, Chapter 8 summarizes the achieved results and outlines possible directions for future improvements.
